\section{Threats to validity}
\label{sec:threats}

\subsection{Reference-set construction, independence and authorship}
\label{sec:threats-gt}
\paragraph{How ‹any evaluated system› was involved.} ‹Facts, not reassurance.›
\paragraph{Possible biases, in both directions.}
\begin{itemize}
  \item \emph{‹Bias 1›.} ‹Mechanism; which direction it could bias the results; measured size or "cannot be tested".›
  \item \emph{‹Bias 2›.} ‹…›
\end{itemize}
‹The directions differ, and the net effect is ‹known / unknown›.›
\paragraph{Authorship.} ‹Employers, products, roles; any author who owns an evaluated system; post-disclosure fixes and that results are as run.›

\subsection{Remaining threats}
\label{sec:threats-other}
\paragraph{Runs per cell.} ‹…›
\paragraph{Reference-set defects.} ‹Mechanism and impact; frozen set used.›
\paragraph{Completeness of the reference set.} ‹Precision as a lower bound; recall can move either way.›
\paragraph{Scoring choices.} ‹…›
\paragraph{Configurations, not models.} ‹…; if only one system of a category was tested, say so and name it wherever a number is attached.›
\paragraph{Execution environment.} ‹…›
\paragraph{Version recording.} ‹What is recorded per run, what is inferred, what is operator-reported.›
\paragraph{Cost.} ‹Known cost only; unavailable is not zero.›
\paragraph{Statistics.} ‹…›
\paragraph{Corpus scope.} ‹…›
